% ============================================================
%  EXAMEN FINAL (SIN RESPUESTAS)
%  Curso de Introducción a la Universidad — Precálculo y Cálculo I
%  Autor: Ing. Gabriel Astudillo
%  Cubre el contenido de las 12 clases del curso (2 horas cada una)
% ============================================================

\documentclass[12pt, a4paper]{article}

% ── Paquetes ─────────────────────────────────────────────────

\usepackage{mathwizards-guia}
\mwguia{Examen Final}{Ing. Gabriel Astudillo $\cdot$ Curso Preuniversitario}

\begin{document}
% ============================================================

% ── Portada / Título ─────────────────────────────────────────
\mwportada{Examen Final}{Examen Final Integral}{Repaso y Evaluación de Todo el Curso: Álgebra, Funciones, Límites, Derivadas y Aplicaciones}

\vspace{4pt}
\begin{cuadroinfo}
  \small
  \textbf{Presentación:} Este examen condensa los temas principales de las 12 clases del curso como una evaluación integradora. Está organizado por bloques correspondientes a cada una de las guías. Te recomiendo resolverlo \textbf{sin mirar las soluciones}, como si fuera un examen real. Cada sección incluye ejercicios resueltos (para repasar la técnica) y un conjunto propuesto (para evaluar tu dominio).
  \\[4pt]
  \textbf{Nivel de Dificultad:}\quad
  \facil{} Básico / Directo \quad
  \medio{} Intermedio \quad
  \dificil{} Desafiante \quad
  \muyDificil{} Muy Difícil / Reto
\end{cuadroinfo}

\vspace{8pt}

% ============================================================
\section{Álgebra, Ecuaciones y Desigualdades}
% ============================================================

\subsection{Ejercicio resuelto integrador}

\textbf{Ejemplo R1.} Resuelve el sistema combinado de ecuaciones y desigualdades:
\begin{enumerate}[label=\roman*)]
  \item $x^2 - 5x + 6 = 0$
  \item $|x - 3| < 2$
\end{enumerate}

\begin{cuadrosolucion}
  \textbf{Solución:}
  \begin{itemize}
    \item Ecuación: $x^2 - 5x + 6 = (x - 2)(x - 3) = 0 \Rightarrow x = 2$ o $x = 3$.
    \item Desigualdad: $-2 < x - 3 < 2 \Rightarrow 1 < x < 5$.
  \end{itemize}
  Las soluciones de la ecuación que cumplen la desigualdad son $x = 2$ y $x = 3$ (ambas están en $(1, 5)$).
\end{cuadrosolucion}

\subsection{Ejercicios propuestos}

\begin{enumerate}[series=ejercicios, leftmargin=*]
  \item \facil{} Simplifica: $\dfrac{(2x^2y)^3 \cdot y^{-2}}{x^{4} y^{2}}$

  \item \facil{} Factoriza completamente: $4x^2 - 12x + 9$.

  \item \facil{} Resuelve: $\dfrac{x + 2}{x - 1} = \dfrac{2x}{x + 3}$.

  \item \medio{} Resuelve la desigualdad: $x^2 - x - 6 \leq 0$.

  \item \medio{} Resuelve: $|2x + 1| \geq 5$.

  \item \dificil{} Racionaliza el numerador de $\dfrac{\sqrt{x + 9} - 3}{x}$ y simplifica.

  \item \dificil{} Resuelve la desigualdad racional: $\dfrac{x^2 - 4}{x - 1} \geq 0$.

  \item \dificil{} Encuentra dos números cuya suma sea 15 y la suma de sus cuadrados sea mínima.
\end{enumerate}

\vspace{4pt}

% ============================================================
\section{Funciones, Dominios y Transformaciones}
% ============================================================

\subsection{Ejercicio resuelto integrador}

\textbf{Ejemplo R1.} Dada $f(x) = x^2 - 2x$, determina: dominio, vértice, interceptos, y la transformación que lleva $y = x^2$ a esta función.

\begin{cuadrosolucion}
  \textbf{Solución:}
  \begin{itemize}
    \item Dominio: $\mathbb{R}$.
    \item Completando el cuadrado: $f(x) = (x - 1)^2 - 1$. Vértice: $(1, -1)$.
    \item Interceptos: $y$: $(0, 0)$; $x$: $x(x - 2) = 0 \Rightarrow x = 0, 2$, así que $(0, 0)$ y $(2, 0)$.
    \item De $y = x^2$ a $y = (x - 1)^2 - 1$: traslación 1 unidad a la derecha y 1 unidad hacia abajo.
  \end{itemize}
\end{cuadrosolucion}

\subsection{Ejercicios propuestos}

\begin{enumerate}[resume=ejercicios, leftmargin=*]
  \item \facil{} Encuentra el dominio de $f(x) = \dfrac{x + 1}{x^2 - 4}$.

  \item \facil{} Encuentra el dominio de $f(x) = \sqrt{2x + 6}$.

  \item \facil{} Dadas $f(x) = x^2$ y $g(x) = x + 2$, calcula $(f \circ g)(x)$ y $(g \circ f)(x)$.

  \item \medio{} Determina si $f(x) = x^2 + 4x$ es inyectiva. Si no lo es, indica una restricción apropiada del dominio.

  \item \medio{} Halla la inversa de $f(x) = \dfrac{2x - 3}{x + 1}$.

  \item \medio{} La gráfica de $f(x) = \sqrt{x}$ se refleja sobre el eje $x$ y se traslada 2 unidades a la derecha y 3 hacia arriba. Escribe la nueva función.

  \item \dificil{} Determina el dominio de $f(x) = \sqrt{\dfrac{x - 1}{x + 2}}$.

  \item \dificil{} Una parábola tiene vértice en $(3, -2)$ y pasa por el punto $(1, 0)$. Encuentra su ecuación.
\end{enumerate}

\newpage

% ============================================================
\section{Polinomiales, Racionales y Composiciones}
% ============================================================

\subsection{Ejercicio resuelto integrador}

\textbf{Ejemplo R1.} Factoriza $f(x) = x^3 - 2x^2 - 5x + 6$ y determina sus asíntotas si es racional.

\begin{cuadrosolucion}
  \textbf{Solución:}
  \begin{enumerate}
    \item Probando $x = 1$: $f(1) = 1 - 2 - 5 + 6 = 0$. Es raíz.
    \item División sintética con 1: cociente $x^2 - x - 6 = (x - 3)(x + 2)$.
    \item Factorización: $f(x) = (x - 1)(x - 3)(x + 2)$.
  \end{enumerate}
  No es racional (es polinomial), así que no tiene asíntotas; su dominio es $\mathbb{R}$.
\end{cuadrosolucion}

\subsection{Ejercicios propuestos}

\begin{enumerate}[resume=ejercicios, leftmargin=*]
  \item \facil{} Divide $(x^3 - 2x^2 + x - 5)$ entre $(x - 2)$ usando división sintética e indica el residuo.

  \item \facil{} Encuentra las asíntotas de $f(x) = \dfrac{2x}{x - 1}$.

  \item \facil{} Dadas $f(x) = \dfrac{1}{x}$ y $g(x) = x + 1$: calcula $(f \circ g)(x)$ y su dominio.

  \item \medio{} Encuentra el dominio, asíntotas, interceptos y huecos de $f(x) = \dfrac{x^2 - 1}{x^2 + x - 2}$.

  \item \medio{} Un polinomio de grado 3 tiene raíces $x = -2$, $x = 1$, $x = 3$ y pasa por $(0, -6)$. Encuentra $f(x)$.

  \item \medio{} Encuentra la inversa de $f(x) = x^3 + 2$.

  \item \dificil{} Encuentra la asíntota oblicua de $f(x) = \dfrac{x^3 + 2x^2 - x + 1}{x^2 + 1}$.

  \item \dificil{} Determina los valores de $k$ tal que $f(x) = \dfrac{kx^2 + 1}{x - 2}$ tenga una asíntota vertical y una asíntota horizontal a la vez.
\end{enumerate}

\newpage

% ============================================================
\section{Exponenciales, Logaritmos y Trigonometría}
% ============================================================

\subsection{Ejercicio resuelto integrador}

\textbf{Ejemplo R1.} Resuelve $\log_2 x + \log_2(x - 2) = 3$ y verifica.

\begin{cuadrosolucion}
  \textbf{Solución:}
  $$\log_2[x(x - 2)] = 3 \Rightarrow x(x - 2) = 8 \Rightarrow x^2 - 2x - 8 = 0 \Rightarrow (x - 4)(x + 2) = 0.$$
  $x = 4$ es válida (argumentos positivos); $x = -2$ se descarta.
\end{cuadrosolucion}

\subsection{Ejercicios propuestos}

\begin{enumerate}[resume=ejercicios, leftmargin=*]
  \item \facil{} Resuelve: $2^x = 32$.

  \item \facil{} Resuelve: $\log_3(x + 1) = 2$.

  \item \facil{} Convierte $135^\circ$ a radianes.

  \item \facil{} Calcula: $\sin\left(\frac{\pi}{3}\right)$ y $\cos\left(\frac{\pi}{6}\right)$.

  \item \medio{} Determina amplitud y período de $y = 3\cos\left(\frac{x}{2}\right)$.

  \item \medio{} Encuentra todos los ángulos en $[0, 2\pi)$ tales que $\cos\theta = \dfrac{\sqrt{2}}{2}$.

  \item \dificil{} Una población crece según $P(t) = 1000e^{0{,}05t}$ con $t$ en años. ¿En cuántos años se duplicará?

  \item \dificil{} Verifica: $\dfrac{\sin^2 x}{\cos x} = \sec x - \cos x$.
\end{enumerate}

\newpage

% ============================================================
\section{Límites y Continuidad}
% ============================================================

\subsection{Ejercicio resuelto integrador}

\textbf{Ejemplo R1.} Calcula $\lim_{x \to 0} \dfrac{\sin(3x)}{x}$ y clasifica la discontinuidad de $f(x) = \dfrac{\sin(3x)}{x}$ en $x = 0$.

\begin{cuadrosolucion}
  \textbf{Solución:}
  $$\lim_{x \to 0} \dfrac{\sin(3x)}{x} = 3 \cdot \lim_{x \to 0} \dfrac{\sin(3x)}{3x} = 3 \cdot 1 = 3.$$
  La función no está definida en $x = 0$ (división entre cero), pero el límite existe. Es una discontinuidad \textbf{evitable}: se puede reparar definiendo $f(0) = 3$.
\end{cuadrosolucion}

\subsection{Ejercicios propuestos}

\begin{enumerate}[resume=ejercicios, leftmargin=*]
  \item \facil{} Calcula: $\lim_{x \to 2} \dfrac{x^2 - 4}{x - 2}$.

  \item \facil{} Calcula: $\lim_{x \to \infty} \dfrac{3x + 1}{x - 5}$.

  \item \facil{} Determina si $f(x) = x^3 - x$ es continua en $x = 1$.

  \item \medio{} Calcula: $\lim_{x \to 0} \dfrac{1 - \cos x}{x^2}$.

  \item \medio{} Encuentra $k$ para que la función sea continua en $x = 2$:
        $$f(x) = \begin{cases} kx + 1 & \text{si } x < 2 \\ 3x - 2 & \text{si } x \geq 2 \end{cases}$$

  \item \medio{} Determina las asíntotas de $f(x) = \dfrac{2x^2 + 3x - 1}{x + 2}$.

  \item \dificil{} Calcula: $\lim_{x \to \infty} \left(\sqrt{x^2 + 1} - x\right)$ (pista: racionaliza).

  \item \dificil{} Usa el TVI para demostrar que $f(x) = x^3 - x - 1$ tiene una raíz en $(1, 2)$.
\end{enumerate}

\newpage

% ============================================================
\section{Derivadas}
% ============================================================

\subsection{Ejercicio resuelto integrador}

\textbf{Ejemplo R1.} Calcula $f'(x)$ para $f(x) = x^2 e^{2x} \sin x$ (aplicando regla del producto y cadena).

\begin{cuadrosolucion}
  \textbf{Solución:} Agrupamos $u = x^2 e^{2x}$, $v = \sin x$:
  \begin{itemize}
    \item $u' = 2x e^{2x} + x^2 \cdot 2e^{2x} = e^{2x}(2x + 2x^2)$.
    \item $v' = \cos x$.
  \end{itemize}
  $$f'(x) = u'v + uv' = e^{2x}(2x + 2x^2)\sin x + x^2 e^{2x} \cos x = e^{2x}\left[(2x + 2x^2)\sin x + x^2 \cos x\right].$$
\end{cuadrosolucion}

\subsection{Ejercicios propuestos}

\begin{enumerate}[resume=ejercicios, leftmargin=*]
  \item \facil{} Calcula la derivada de $f(x) = 3x^4 - 2x^3 + x^2 - 5$.

  \item \facil{} Calcula la derivada de $f(x) = e^x + \ln x + \sin x$.

  \item \facil{} Calcula la derivada de $f(x) = (x^2 + 3x)^5$.

  \item \facil{} Calcula $f''(x)$ para $f(x) = e^{2x}$.

  \item \medio{} Calcula la derivada de $f(x) = \dfrac{x^2}{x^2 + 1}$.

  \item \medio{} Calcula la derivada de $f(x) = \ln(x^2 + 1)$.

  \item \medio{} Usando la definición de derivada, calcula $f'(x)$ para $f(x) = \dfrac{1}{x + 1}$.

  \item \dificil{} Calcula la derivada de $f(x) = \sqrt{\sin x + e^x}$.

  \item \dificil{} Encuentra la ecuación de la recta tangente a $f(x) = x^3$ en $x = 2$.

  \item \dificil{} Determina los valores de $x$ donde la recta tangente a $f(x) = x^3 - 6x^2 + 5$ es horizontal.
\end{enumerate}

\newpage

% ============================================================
\section{Aplicaciones de la Derivada}
% ============================================================

\subsection{Ejercicio resuelto integrador}

\textbf{Ejemplo R1.} Encuentra los extremos absolutos de $f(x) = x^3 - 3x^2 + 1$ en $[-1, 4]$.

\begin{cuadrosolucion}
  \textbf{Solución:}
  \begin{itemize}
    \item $f'(x) = 3x^2 - 6x = 3x(x - 2)$. Puntos críticos: $x = 0$, $x = 2$ (ambos en el intervalo).
    \item Evaluamos: $f(-1) = -1 - 3 + 1 = -3$; $f(0) = 1$; $f(2) = 8 - 12 + 1 = -3$; $f(4) = 64 - 48 + 1 = 17$.
  \end{itemize}
  Máximo absoluto: $17$ en $x = 4$. Mínimo absoluto: $-3$ (en $x = -1$ y $x = 2$).
\end{cuadrosolucion}

\subsection{Ejercicios propuestos}

\begin{enumerate}[resume=ejercicios, leftmargin=*]
  \item \facil{} Encuentra los puntos críticos de $f(x) = x^2 - 8x + 12$.

  \item \facil{} Determina los intervalos de crecimiento de $f(x) = x^2 - 4x$.

  \item \facil{} Aplica el TVM a $f(x) = x^2$ en $[1, 3]$ y encuentra $c$.

  \item \medio{} Encuentra las dimensiones del rectángulo de área máxima con perímetro $24$ m.

  \item \medio{} Se lanza una pelota con altura $h(t) = -5t^2 + 20t + 2$ metros. ¿Cuál es la altura máxima?

  \item \medio{} Determina intervalos de concavidad de $f(x) = x^3 - 3x^2$.

  \item \dificil{} Un tanque cilíndrico con tapa tiene volumen fijo $V = 1000$ cm$^3$. Encuentra las dimensiones que minimizan el material (área total).

  \item \dificil{} Una escalera de $5$ m está apoyada contra una pared. La base se desliza a $1$ m/s. ¿Qué tan rápido cae la parte superior cuando la base está a $3$ m de la pared?

  \item \dificil{} Encuentra los extremos absolutos de $f(x) = x \ln x$ en $[1, e]$.

  \item \dificil{} Determina los puntos de inflexión de $f(x) = x^4 - 6x^2 + 4x$.
\end{enumerate}
\end{document}